% !TEX program = pdflatex
\documentclass[11pt,a4paper]{article}

\usepackage[margin=2cm]{geometry}
\usepackage[T1]{fontenc}
\usepackage[utf8]{inputenc}
\usepackage[spanish]{babel}
\usepackage{lmodern}
\usepackage{amsmath,amssymb,amsthm}
\usepackage{mathrsfs}
\usepackage{titlesec}
\usepackage{microtype}
\usepackage{fancyhdr}
\usepackage[hidelinks]{hyperref}

\titleformat{\section}{\Large\bfseries}{\thesection}{0.7em}{}
\titlespacing*{\section}{0pt}{1.5em}{0.6em}
\setlength{\parindent}{0pt}
\setlength{\parskip}{0.55em}

\newtheorem{teorema}{Teorema}
\newtheorem{observacion}{Observación}

\setlength{\headheight}{14pt}
\setlength{\headsep}{0.5cm}
\pagestyle{fancy}
\fancyhf{}
\fancyhead[L]{\small Grupo de Estudio: Teoría Analítica de Números}
\fancyhead[R]{\small\today}
\renewcommand{\headrulewidth}{0.4pt}
\fancypagestyle{plain}{%
  \fancyhf{}
  \fancyhead[L]{\small Grupo de Estudio: Teoría Analítica de Números}
  \fancyhead[R]{\small\today}
  \renewcommand{\headrulewidth}{0.4pt}
}
\newcommand{\N}{\mathbb{N}}

\title{Actividades de la Clase 0}
\author{}
\date{}

\begin{document}
\maketitle
\noindent{\small Autor: Daniel Soto Villada}\par


La factorización única de los enteros permite traducir una suma sobre todos los
números naturales a un producto sobre los primos. La prueba se construye primero
para un conjunto finito de primos.

En todo el documento, usamos $\N=\{1,2,3,\ldots\}$ y
$\mathbb{N}_0=\{0,1,2,\ldots\}$.

\section{Enunciado}

Sea $s\in\mathbb{C}$ y escribamos $\sigma=\Re(s)$. Para
$\sigma>1$, la serie
\[
  \zeta(s)=\sum_{n=1}^{\infty}\frac{1}{n^s}
\]
converge absolutamente. En esta región probaremos que
\[
  \zeta(s)=\prod_{p\ \mathrm{primo}}\frac{1}{1-p^{-s}}
  \qquad (\Re(s)>1).
\]
El producto infinito se entenderá como el límite de los productos finitos
tomados sobre los primos $p\le y$.

\section{Demostración}

Para un primo $p$ y un entero $n\in\mathbb{N}$, definimos la valoración p-ádica
\[
  v_p(n):=\max\{a\in\mathbb{N}_0:p^a\mid n\}.
\]
Es decir, $v_p(n)$ es el mayor exponente para el cual $p^a$ todavía divide a
$n$.

Para $n=1$, se tiene $v_p(1)=0$ para todo primo $p$. Por el teorema
fundamental de la aritmética, todo entero $n\ge2$ se escribe de manera única
como producto de potencias primas; para $n=1$ usamos el producto vacío. Así,
\[
  n=\prod_{p\le n}p^{v_p(n)},\qquad
  n^{-s}=\prod_{p\le n}p^{-v_p(n)s}.
\]
Al sustituir esta expresión en la serie de zeta, obtenemos
\[
  \zeta(s)=\sum_{n=1}^{\infty}\prod_{p\le n}p^{-v_p(n)s}.
\]
La unicidad de la factorización dice que cada entero aparece exactamente una vez
cuando combinamos potencias de primos. Veamos cómo esa combinación produce un
producto. Primero usamos solo el primo $2$:
\[
  1+2^{-s}+2^{-2s}+2^{-3s}+\cdots=\frac{1}{1-2^{-s}}.
\]
Al incorporar también el primo $3$, aplicamos la distributividad y agrupamos
los términos según la potencia de $2$:
\begin{align*}
 &(1+2^{-s}+2^{-2s}+2^{-3s}+\cdots)
   (1+3^{-s}+3^{-2s}+3^{-3s}+\cdots)\\
 &=\underbrace{(1+3^{-s}+3^{-2s}+3^{-3s}+\cdots)}_{\text{potencia }2^0}\\
 &\quad+\underbrace{(2^{-s}+2^{-s}3^{-s}+2^{-s}3^{-2s}+2^{-s}3^{-3s}+\cdots)}_{\text{potencia }2^1}\\
 &\quad+\underbrace{(2^{-2s}+2^{-2s}3^{-s}+2^{-2s}3^{-2s}+2^{-2s}3^{-3s}+\cdots)}_{\text{potencia }2^2}\\
 &\quad+\underbrace{(2^{-3s}+2^{-3s}3^{-s}+2^{-3s}3^{-2s}+2^{-3s}3^{-3s}+\cdots)}_{\text{potencia }2^3}+\cdots\\
 &=\sum_{a,b\ge0}2^{-as}3^{-bs}
  =\sum_{a,b\ge0}\frac{1}{(2^a3^b)^s}.
\end{align*}
Las primeras filas aportan respectivamente los denominadores
\[
  1,3,9,27,\ldots;\qquad 2,6,18,54,\ldots;\qquad
  4,12,36,108,\ldots;\qquad 8,24,72,\ldots.
\]
Así aparecen una vez cada uno los enteros cuyos únicos factores primos son
$2$ y $3$. Por ejemplo, $12=2^2\cdot3$ aparece en la fila de $2^2$ y en la
columna de $3^1$. Incorporar $5$, luego $7$ y después los demás primos repite
el mismo procedimiento.

Fijemos un entero $y\ge2$. Denotemos por
\[
  \mathbb{P}_y=\{p\le y:p\text{ es primo}\},\qquad
  \mathscr{P}_y(s)=\prod_{p\in\mathbb{P}_y}(1-p^{-s})^{-1}.
\]
Sea también
\[
  \mathbb{S}_y=\{n\in\mathbb{N}: v_p(n)>0\Rightarrow p\le y
  \text{ para todo primo }p\}.
\]
En particular, $1\in\mathbb{S}_y$, pues $v_p(1)=0$ para todo primo $p$.

Por ejemplo, para $y=5$ los primos permitidos son $2$, $3$ y $5$, y
\[
  \mathbb{S}_5=\{1,2,3,4,5,6,8,9,10,12,15,16,18,20,24,25,27,30,\ldots\}.
\]
En efecto, $12=2^2\cdot3$, $25=5^2$ y $30=2\cdot3\cdot5$ pertenecen a
$\mathbb{S}_5$. En cambio, $14=2\cdot7$ no pertenece, porque contiene el
factor primo $7>5$.

Expandimos cada factor. Para un primo fijo $p\le y$, como
$|p^{-s}|=p^{-\sigma}<1$, la serie geométrica converge y
\[
  (1-p^{-s})^{-1}=\sum_{k=0}^{\infty}p^{-ks}
  =1+p^{-s}+p^{-2s}+p^{-3s}+\cdots.
\]
Sustituyendo esta igualdad para cada primo de $\mathbb{P}_y$ en el producto
finito, tenemos
\[
  \mathscr{P}_y(s)=\prod_{p\in\mathbb{P}_y}
          \left(\sum_{k_p=0}^{\infty}p^{-k_ps}\right).
\]
Veamos la distribución completa para $y=5$, donde los primos permitidos son
$2$, $3$ y $5$. Hay tres paréntesis, y se elige un término de cada uno:
\begin{align*}
 P_5(s)
 &=\left(\sum_{k_2=0}^{\infty}2^{-k_2s}\right)
   \left(\sum_{k_3=0}^{\infty}3^{-k_3s}\right)
   \left(\sum_{k_5=0}^{\infty}5^{-k_5s}\right)\\
 &=\sum_{k_2=0}^{\infty}\sum_{k_3=0}^{\infty}\sum_{k_5=0}^{\infty}
   2^{-k_2s}3^{-k_3s}5^{-k_5s}.
\end{align*}
Esto es la propiedad distributiva de siempre, aplicada a tres sumas: cada
combinación de exponentes $(k_2,k_3,k_5)$ aporta un término. Por ejemplo,
$(2,1,0)$ aporta $2^{-2s}3^{-s}5^0=12^{-s}$, mientras que $(0,0,2)$ aporta
$25^{-s}$.

Para un valor general de $y$, hacemos lo mismo: elegimos un exponente no
negativo para cada primo $p\in\mathbb{P}_y$. En forma abreviada,
\[
  \mathscr{P}_y(s)=\sum_{(k_p)\in\mathbb{N}_0^{\mathbb{P}_y}}
          \prod_{p\in\mathbb{P}_y}p^{-k_ps}.
\]
Aquí $\mathbb{N}_0^{\mathbb{P}_y}$ significa el conjunto de todas las listas
$(k_p)_{p\in\mathbb{P}_y}$: una lista contiene un exponente para cada primo
permitido.

La distribución y la reagrupación son válidas porque las series convergen
absolutamente. En efecto, la suma de los valores absolutos de los términos es
\[
  \sum_{(k_p)\in\mathbb{N}_0^{\mathbb{P}_y}}
     \left|\prod_{p\in\mathbb{P}_y}p^{-k_ps}\right|
  =\prod_{p\in\mathbb{P}_y}\left(\sum_{k=0}^{\infty}p^{-k\sigma}\right)
  =\prod_{p\in\mathbb{P}_y}(1-p^{-\sigma})^{-1}<\infty.
\]

Cada lista de exponentes determina un entero.A la lista
$(k_p)_{p\in\mathbb{P}_y}$ le asociamos el número
\[
  n=\prod_{p\in\mathbb{P}_y}p^{k_p}.
\]
Veamos cómo se calcula este producto con números concretos. Si $y=5$, los
primos permitidos son $2$, $3$ y $5$, así que la fórmula queda
$n=2^{k_2}3^{k_3}5^{k_5}$. Por ejemplo, para
$(k_2,k_3,k_5)=(2,1,0)$ obtenemos
$n=2^2\cdot3^1\cdot5^0=4\cdot3\cdot1=12$. Si elegimos
$(k_2,k_3,k_5)=(0,2,1)$, obtenemos
$n=2^0\cdot3^2\cdot5^1=1\cdot9\cdot5=45$. Cada lista de exponentes
produce así un número $n$.

Por definición, $\mathbb{S}_y$ contiene los números que no tienen
ningún factor primo mayor que $y$. En la fórmula de $n$ solo multiplicamos
primos $p\in\mathbb{P}_y$, es decir, primos que cumplen $p\le y$. Elevar
uno de esos primos a una potencia no introduce factores primos nuevos. Por
tanto, todos los factores primos de $n$ son $\le y$, y eso es precisamente
lo que significa que $n\in\mathbb{S}_y$.

En el primer ejemplo, los únicos factores primos de $12$ son $2$ y $3$,
ambos menores o iguales que $5$; por eso $12\in\mathbb{S}_5$. Además, las
leyes de los exponentes muestran que el sumando correspondiente a cualquier
lista es precisamente
\[
  \prod_{p\in\mathbb{P}_y}p^{-k_ps}
  =\left(\prod_{p\in\mathbb{P}_y}p^{k_p}\right)^{-s}
  =n^{-s}.
\]

El teorema fundamental garantiza que no falta ni se repite ningún
entero. Recíprocamente, si $n\in\mathbb{S}_y$, todos sus factores primos
pertenecen a $\mathbb{P}_y$. Su factorización prima única proporciona una
lista de exponentes $(k_p)_{p\in\mathbb{P}_y}$: usamos $k_p=0$ para los primos
permitidos que no dividen a $n$. La unicidad de la factorización garantiza que
esa lista es la única que produce $n$. Por tanto, las listas de exponentes
corresponden exactamente, una a una, a los enteros de $\mathbb{S}_y$, y
\[
  \sum_{n\in\mathbb{S}_y}n^{-s}
  =\mathscr{P}_y(s)
  =\prod_{p\in\mathbb{P}_y}(1-p^{-s})^{-1},
  \qquad y\ge2,\ \Re(s)>1.
\]

Hasta aquí queda demostrado el caso de un conjunto finito de primos: el producto
$\mathscr{P}_y(s)$ tiene finitos factores, aunque cada factor geométrico sea una serie
infinita. Para obtener toda la serie $\zeta(s)$ todavía debemos hacer tender
a $y$ al infinito; ese paso se justifica más adelante.

\end{document}